
\documentclass{article}
\usepackage{colortbl}
\usepackage{makecell}
\usepackage{multirow}
\usepackage{supertabular}

\begin{document}

\newcounter{utterance}

\centering \large Interaction Transcript for game `eqbench', experiment `eq\_bench', episode 70 with reflexion{-}Qwen3.5{-}9B.
\vspace{24pt}

{ \footnotesize  \setcounter{utterance}{1}
\setlength{\tabcolsep}{0pt}
\begin{supertabular}{c@{$\;$}|p{.15\linewidth}@{}p{.15\linewidth}p{.15\linewidth}p{.15\linewidth}p{.15\linewidth}p{.15\linewidth}}
    \# & \multicolumn{2}{c}{Player} && \multicolumn{2}{c}{Game Master} \\
    \hline

    \theutterance \stepcounter{utterance}  
    & & & \multicolumn{4}{p{0.6\linewidth}}{
        \cellcolor[rgb]{0.9,0.9,0.9}{
            \makecell[{{p{\linewidth}}}]{
                \texttt{\tiny{[P1$\langle$GM]}}
                \texttt{You have attempted this same task before and did not succeed. These are the plans you wrote after those attempts.} \\
\texttt{Trial \#1: New plan: I failed because I assigned a high score to "Resentment" (9) and a low score to "Pride" (2), missing the core dynamic of the dialogue. Jacques is a chef who is being insulted about his cooking skills and ego; his aggressive, defensive, and boastful responses ("more talent in my little finger," "disgrace to the culinary world") indicate that his primary emotional driver is wounded **Pride**, not just general resentment. While he is angry, the specific nature of the insults targets his self{-}worth and professional identity, making Pride the dominant emotion. Next time, I will analyze the specific content of the insults (ego, talent, career) to identify that the character's defensive aggression is a direct manifestation of bruised Pride, and I will assign the highest score (10) to Pride, with Resentment as a secondary, lower score (e.g., 4), and zero for Fear and Apathy.} \\
\\ 
\texttt{Trial \#2: New plan: I failed because I misidentified the primary emotional driver as "Resentment" (4) while underestimating "Pride" (10). The dialogue shows Jacques reacting defensively to insults about his cooking skills and ego ("more talent in my little finger," "disgrace to the culinary world"). His aggression is a direct defense of his wounded self{-}worth and professional identity, not just passive bitterness. Next time, I will recognize that the specific content of the insults targets his ego and career, making **Pride** the dominant emotion (score 10). I will assign a secondary score to **Resentment** (e.g., 4) to reflect his anger, and keep **Fear** and **Apathy** at 0, as there is no evidence of fear or indifference in his responses.} \\
\\ 
\\ 
\\ 
\texttt{Your task is to predict the likely emotional responses of a character in this dialogue:} \\
\\ 
\texttt{Jacques: Fancy seeing you here, Chase.} \\
\texttt{Olivia: The feeling is mutual, Duval. } \\
\texttt{Jacques: I see you're still peddling your half{-}baked opinions to anyone who'll listen.} \\
\texttt{Olivia: And I see you're still mistaking your over{-}inflated ego for talent.} \\
\texttt{Jacques: Talent? I've got more talent in my little finger than you have in your entire body.} \\
\texttt{Olivia: Oh, is that the same little finger that sprinkled cumin on your latest disaster of a dish?} \\
\texttt{Jacques: Cumin? Are you out of your mind, woman? It was caraway.} \\
\texttt{Olivia: Caraway? Don't make me laugh. I know cumin when I taste it.} \\
\texttt{Jacques: You wouldn't know cumin if it bit you on the ass.} \\
\texttt{Olivia: I'd certainly know it better than you, it seems.} \\
\texttt{Jacques: You have no idea what you're talking about. You're a hack, Chase.} \\
\texttt{Olivia: And you're a has{-}been, Duval. } \\
\texttt{Jacques: That's it. I've had enough of your insolence.} \\
\texttt{Olivia: And I've had enough of your mediocrity.} \\
\texttt{Jacques: You're a disgrace to the culinary world.} \\
\texttt{Olivia: And you're a disgrace to chefs everywhere.} \\
\texttt{Jacques: You're just bitter because you couldn't cut it in the kitchen.} \\
\texttt{Olivia: And you're just angry because your star is fading.} \\
\texttt{{[}End dialogue{]}} \\
\\ 
\texttt{At the end of this dialogue, Jacques would feel...} \\
\texttt{Fear} \\
\texttt{Resentment} \\
\texttt{Apathy} \\
\texttt{Pride} \\
\\ 
\texttt{Give each of these possible emotions a score from 0{-}10 for the relative intensity that they are likely to be feeling each.} \\
\\ 
\texttt{You must output in the following format, including headings (of course, you should give your own scores), with no additional commentary:} \\
\\ 
\texttt{Fear: <score>} \\
\texttt{Resentment: <score>} \\
\texttt{Apathy: <score>} \\
\texttt{Pride: <score>} \\
\\ 
\\ 
\texttt{{[}End of answer{]}} \\
\\ 
\texttt{Remember: zero is a valid score, meaning they are likely not feeling that emotion. You must score at least one emotion > 0.} \\
\\ 
\texttt{Your answer:} \\
            }
        }
    }
     \\ \\

    \theutterance \stepcounter{utterance}  
    & \multicolumn{4}{p{0.6\linewidth}}{
        \cellcolor[rgb]{0.561,0.737,0.561}{
            \makecell[{{p{\linewidth}}}]{
                \texttt{\tiny{[P1$\rangle$GM]}}
                \texttt{Fear: 0} \\
\texttt{Resentment: 4} \\
\texttt{Apathy: 0} \\
\texttt{Pride: 10} \\
            }
        }
    }
    & & \\ \\

    \theutterance \stepcounter{utterance}  
    & & & \multicolumn{2}{p{0.3\linewidth}}{
        \cellcolor[rgb]{0.9,0.9,0.9}{
            \makecell[{{p{\linewidth}}}]{
                \texttt{\tiny{[GM$|$GM]}}
                \texttt{Fear: 0} \\
\texttt{Resentment: 4} \\
\texttt{Apathy: 0} \\
\texttt{Pride: 10} \\
            }
        }
    }
    & & \\ \\

    \theutterance \stepcounter{utterance}  
    & & & \multicolumn{2}{p{0.3\linewidth}}{
        \cellcolor[rgb]{0.9,0.9,0.9}{
            \makecell[{{p{\linewidth}}}]{
                \texttt{\tiny{[GM$|$GM]}}
                \texttt{Fear: 0} \\
\texttt{Resentment: 8} \\
\texttt{Apathy: 0} \\
\texttt{Pride: 4} \\
            }
        }
    }
    & & \\ \\

    \theutterance \stepcounter{utterance}  
    & & & \multicolumn{2}{p{0.3\linewidth}}{
        \cellcolor[rgb]{0.9,0.9,0.9}{
            \makecell[{{p{\linewidth}}}]{
                \texttt{\tiny{[GM$|$GM]}}
                \texttt{game\_result = LOSE} \\
            }
        }
    }
    & & \\ \\

    \theutterance \stepcounter{utterance}  
    & & \multicolumn{4}{p{0.6\linewidth}}{
        \cellcolor[rgb]{0.95,0.95,0.95}{
            \makecell[{{p{\linewidth}}}]{
                \texttt{\tiny{[P1$\leftrightarrow$P]}}
                \texttt{Post{-}episode Reflexion (not sent to GM):} \\
\texttt{New plan: I failed because I misidentified the primary emotional driver as "Resentment" (4) while underestimating "Pride" (10). The dialogue shows Jacques reacting defensively to insults about his cooking skills and ego ("more talent in my little finger," "disgrace to the culinary world"). His aggression is a direct defense of his wounded self{-}worth and professional identity, not just passive bitterness. Next time, I will recognize that the specific content of the insults targets his ego and career, making **Pride** the dominant emotion (score 10). I will assign a secondary score to **Resentment** (e.g., 4) to reflect his anger, and keep **Fear** and **Apathy** at 0, as there is no evidence of fear or indifference in his responses.} \\
            }
        }
    }
    & \\ \\

\end{supertabular}
}

\end{document}
